\documentclass[11pt]{article}% Shared preamble for the rigor documents (Lamport-style structured proofs).  \input this file after \documentclass.
\usepackage[T1]{fontenc}\usepackage{lmodern}\usepackage{microtype}
\usepackage[margin=1in]{geometry}
\usepackage{amsmath,amssymb,amsthm,mathtools}
\usepackage{xcolor}\usepackage{enumitem}\usepackage{booktabs}\usepackage{graphicx}\usepackage{hyperref}
\usepackage[most]{tcolorbox}
\definecolor{navy}{HTML}{17365D}\definecolor{teal}{HTML}{117A8B}\definecolor{warn}{HTML}{A33A2B}\definecolor{okgreen}{HTML}{2E7D4F}
\hypersetup{colorlinks=true,linkcolor=navy,citecolor=teal,urlcolor=teal}
\theoremstyle{plain}
\newtheorem{theorem}{Theorem}[section]\newtheorem{lemma}[theorem]{Lemma}\newtheorem{proposition}[theorem]{Proposition}\newtheorem{corollary}[theorem]{Corollary}
\theoremstyle{definition}\newtheorem{definition}[theorem]{Definition}\newtheorem{remark}[theorem]{Remark}\newtheorem{claim}[theorem]{Claim}\newtheorem{issue}[theorem]{Issue}
% ---------- Lamport structured proofs ----------
% \lstep{level}{number}{assertion}   -> indented "<level>number. assertion"
% \lproof                             -> "PROOF:" line (start of the sub-proof of the preceding step)
% \lby{justification}                 -> "BY ..." leaf justification for the preceding step
% \lqed{level}{number}                -> QED step of a sub-proof
% keywords: \lassume \lprove \lsuffices \lcase \ldefine \lpick \lnew
\newlength{\lindent}\setlength{\lindent}{1.7em}
\newcommand{\lstep}[3]{\par\smallskip\noindent\hspace*{\dimexpr\numexpr#1-1\relax\lindent\relax}{\color{navy}$\langle#1\rangle#2$.}\ #3}
\newcommand{\lproof}{\par\noindent\hspace*{\lindent}{\scshape Proof:}}
\newcommand{\lby}[1]{\par\noindent\hspace*{\lindent}{\scshape By}\ #1.}
\newcommand{\lqed}[2]{\par\smallskip\noindent\hspace*{\dimexpr\numexpr#1-1\relax\lindent\relax}{\color{navy}$\langle#1\rangle#2$.}\ {\scshape Q.E.D.}}
\newcommand{\lassume}{{\scshape Assume}\ }\newcommand{\lprove}{{\scshape Prove}\ }\newcommand{\lsuffices}{{\scshape Suffices}\ }
\newcommand{\lcase}{{\scshape Case}\ }\newcommand{\ldefine}{{\scshape Define}\ }\newcommand{\lpick}{{\scshape Pick}\ }\newcommand{\lnew}{{\scshape New}\ }
% ---------- verdict boxes ----------
\newtcolorbox{verdictok}{colback=okgreen!8,colframe=okgreen,title={\bfseries Verdict: verified},fonttitle=\small}
\newtcolorbox{verdictgap}{colback=orange!10,colframe=orange!80!black,title={\bfseries Verdict: gap / caveat},fonttitle=\small}
\newtcolorbox{verdictbreak}{colback=warn!10,colframe=warn,title={\bfseries Verdict: proof-breaking issue},fonttitle=\small}
\newtcolorbox{citedbox}[1]{colback=navy!5,colframe=navy!60,title={\bfseries Cited result: #1},fonttitle=\small,breakable}
\setlength{\parindent}{0pt}\setlength{\parskip}{4pt}\allowdisplaybreaks
\newcommand{\cA}{\mathcal A}\newcommand{\cF}{\mathcal F}\newcommand{\cM}{\mathfrak M}\newcommand{\cN}{\mathfrak N}

% extra calligraphic shortcuts (\providecommand: no clash if a document defines its own)
\providecommand{\cB}{\mathcal B}\providecommand{\cC}{\mathcal C}\providecommand{\cD}{\mathcal D}\providecommand{\cE}{\mathcal E}\providecommand{\cG}{\mathcal G}\providecommand{\cH}{\mathcal H}\providecommand{\cI}{\mathcal I}\providecommand{\cJ}{\mathcal J}\providecommand{\cK}{\mathcal K}\providecommand{\cL}{\mathcal L}\providecommand{\cO}{\mathcal O}\providecommand{\cP}{\mathcal P}\providecommand{\cQ}{\mathcal Q}\providecommand{\cR}{\mathcal R}\providecommand{\cS}{\mathcal S}\providecommand{\cT}{\mathcal T}\providecommand{\cU}{\mathcal U}\providecommand{\cV}{\mathcal V}\providecommand{\cW}{\mathcal W}\providecommand{\cX}{\mathcal X}\providecommand{\cY}{\mathcal Y}\providecommand{\cZ}{\mathcal Z}
\newcommand{\Fid}{\operatorname{F}}\newcommand{\Tr}{\operatorname{Tr}}\newcommand{\eps}{\varepsilon}\newcommand{\dd}{\,\mathrm d}

% --- excerpts of cited sources ----------------------------------------------------
% The page images in rigor/shots/ are excerpts of third-party publications and are NOT
% redistributed with the public release.  \shot is an exact drop-in for \includegraphics
% that degrades to a framed note when shots/ is absent, so every document still compiles
% from a clone without them; \shotc is the centred variant.  The precise locator of each
% excerpt is always stated in the accompanying cited-result box.
\newcommand{\shotmissing}[2]{\fbox{\parbox{\dimexpr#1-2\fboxsep-2\fboxrule\relax}{\small%
\textbf{Screenshot not included.}\ \texttt{\detokenize{#2}} is an image of a third-party
publication, omitted from the public release; see the locator in the cited-result box.}}}
\newcommand{\shot}[2][\textwidth]{%
  \IfFileExists{shots/#2}{\includegraphics[width=#1]{shots/#2}}{\shotmissing{#1}{#2}}}
\newcommand{\shotc}[2][\textwidth]{\begin{center}\shot[#1]{#2}\end{center}}

\usepackage{tabularx}
\newcommand{\norm}[1]{\lVert#1\rVert}\newcommand{\abs}[1]{\lvert#1\rvert}\newcommand{\Gh}{\widehat G}
\newcommand{\Pp}{P_+}\newcommand{\Pm}{P_-}\newcommand{\Stwo}{\mathfrak S_2}\newcommand{\Dc}{\mathfrak D}
\newcommand{\Dtot}{D_{\rm tot}}\newcommand{\Om}{\Omega}\newcommand{\ip}[2]{\langle#1,#2\rangle}
\newcommand{\doc}{\textbf{NSO}}
\DeclareMathOperator{\artanh}{artanh}\DeclareMathOperator{\dist}{dist}\newcommand{\Qt}{\widetilde Q}\newcommand{\one}{\mathbf 1}\newcommand{\Sch}{\mathcal S}\newcommand{\Rp}{\mathsf R}\newcommand{\kt}{\widetilde k}\newcommand{\Wt}{\widetilde W}

\title{\color{navy}\bfseries Referee report REF-P6-4 on \texttt{rigor/network\_second\_order.tex}\\[0.3em]
\large The second-order network law, the corner kernel $\hat G$, and the six G1b cards}
\author{Referee REF-P6-4 (Opus)}\date{2026-10-03}

\begin{document}\maketitle

\begin{abstract}\noindent
Adversarial review of \texttt{rigor/network\_second\_order.tex} (G1b, 27 pp.) and of the six cards built on it.
Every step of T1, T2, T4, T5, T6 and of the lemmas behind them was checked against the cited documents (QL, PA, RR,
KI, TB, C11, G1a; the labels and corrections quoted below were opened, not remembered), and every number in a
theorem or card Statement was recomputed with my own scripts (\texttt{rigor/ref\_p6\_4\_kernel.py},
\texttt{rigor/ref\_p6\_4\_composite.py}; no code shared with \texttt{g1b\_kernel.py}).
\textbf{Verdict.} T1, T2, the network law at fixed configuration (cor:law), the kernel theorem T4 with its closed
form, bounds, sign and monotonicity, T5, cor:vwz and T6 \textbf{stand}.  \textbf{T3 stands only with a repair}: its
range ``every $s$ with $2\hat f_2s^2Z^2<1$'' uses RR Lemma~6.1 (\texttt{lem:groupbound}) ``for all real $\sigma$'',
which RR's own Correction~C3 (l.~1062--1079) settles only for $|\sigma|\,\norm{[P_-,K]}_2<1$; the proved range is
$2\hat f_2s^2Z^2<1$ \emph{and} $s\max_k\zeta_k<\pi\sqrt3$, the rest rests on the conjectural card
\texttt{group-index-zero-large-s}.  Minor points: the same unrestricted ``for $\sigma\ge0$'' in lem:onegroup(i)
(harmless for T5, which needs only $\sigma\to0$); ``$C^1$ but not $C^2$'' is proved only in a remark; a QAWF/QAWO
mislabel in the script docstring.  All independent recomputations agree with the document (largest deviation
$7\cdot10^{-16}$ against tolerances $10^{-8}$--$10^{-12}$); the rerun of \texttt{g1b\_kernel.py} reproduces
\texttt{g1b\_kernel.out} byte for byte.  Card verdicts: five pass, \texttt{network-global-bound-coincident} fails
(repair: restrict the range).
\end{abstract}

\tableofcontents

\section{Method and verdict table}\label{sec:method}
\doc\ is \texttt{rigor/network\_second\_order.tex} as of 2026-10-03 17:51 (sha of its output
\texttt{g1b\_kernel.out}: \texttt{7581c0e3\dots}; unchanged by my rerun).  Its log: 27 pages, no line \verb|^!|, no
undefined reference.  Labels of the cited documents were opened at the lines named here: QL
\texttt{quadratic\_limit.tex} (\texttt{lem:standing} l.~135, \texttt{def:PhiP} l.~212, Bures box l.~266--290,
\texttt{prop:oneparticle} l.~329, \texttt{thm:lower} l.~410--470); PA \texttt{uhlmann\_upper\_bound.tex}
(\texttt{eq:hardy} l.~42, \texttt{lem:kernel} l.~220, \texttt{prop:HS} l.~284, \texttt{def:admissible},
\texttt{prop:expansion} l.~745--800, \texttt{lem:tangentvec} l.~1096, \texttt{prop:infimum} l.~1216, Corrections
C2--C3 l.~1562--1612); RR \texttt{rate\_of\_remainder.tex} (\texttt{thm:global} l.~180, \texttt{lem:groupbound}
l.~818, \texttt{prop:rrr} l.~848, Correction C3 l.~1026--1079); TB \texttt{universality\_theorem\_B.tex} ((BW)
l.~83--89, (H1)--(H3) l.~118--160, \texttt{prop:H3} l.~290, \texttt{rem:H2H3} l.~358, \texttt{prop:H2} l.~392,
\texttt{lem:reduce} l.~517, \texttt{prop:reduce} l.~569, \texttt{thm:linear} l.~594); C11 \texttt{thm:main} l.~792 and
\S8(b) l.~824; G1a \texttt{network\_corner\_calculus.tex} (\texttt{def:H} l.~614, \texttt{thm:corner} l.~658,
\texttt{thm:state} l.~870, \texttt{cor:phiA} l.~924, \texttt{prop:1B1A} l.~1156, \texttt{prop:prot2} l.~1199,
\texttt{thm:sep} l.~1606); \texttt{f3\_t0.py} and \texttt{kkt\_continuum.Qmat}.  No external PDF is cited directly by
\doc; its seven excerpts (\texttt{rigor/shots/g1b\_*.png}) are of project documents and exist.
Numbers: \S\ref{sec:num}.  Verdict words: STANDS; STANDS WITH REPAIR (= \texttt{kb verdict fail} with the repair);
FAILS.

\begin{center}\small
\begin{tabularx}{\textwidth}{>{\raggedright\arraybackslash}p{0.27\textwidth} >{\raggedright\arraybackslash}X >{\raggedright\arraybackslash}p{0.17\textwidth}}\toprule
Statement of \doc & Finding (section) & Verdict\\\midrule
thm:T3 (Thm.~5.2) & range uses RR lem:groupbound beyond RR Correction C3; proved for $s\max_k\zeta_k<\pi\sqrt3$ only (\S\ref{sec:i9}) & \textbf{stands with repair}\\
lem:onegroup(i) (Lem.~7.1) & ``for $\sigma\ge0$'' has the same C3 restriction; only $\sigma\to0$ is used (\S\ref{sec:i12}) & stands (minor)\\
lem:frame, lem:canon, rem:scaled & (H) stated and used through G1a thm:corner (\S\ref{sec:i1}) & stands\\
lem:implementer & every input checked; normality from PA, not from Note 4 (\S\ref{sec:i2}) & stands\\
lem:Dtot, lem:polar & sign and frame $y_k$ confirmed independently (\S\ref{sec:i3}) & stands\\
lem:taylor, thm:T1 & first-order differentiability suffices (\S\S\ref{sec:i4},\ref{sec:i5}) & stands\\
lem:prodexp, prop:inf, thm:T2 & (\S\ref{sec:i6}) & stands\\
cor:law & fixed configuration only, nothing uniform claimed (\S\ref{sec:i7}) & stands\\
lem:nogroup & masses recomputed exactly (\S\ref{sec:i8}) & stands\\
thm:T4 (a)--(e), (cf), lem:lorentz & every step and every number recomputed; ``not $C^2$'' only in rem:kernelshape (\S\ref{sec:i10}) & stands (minor)\\
cor:W3W4 & (\S\ref{sec:i11}) & stands\\
thm:T5, cor:vwz & uniform in $y_2-y_1$; $M=2$ used in $\langle2\rangle1$; strengths recomputed exactly (\S\S\ref{sec:i12},\ref{sec:i13}) & stands\\
thm:T6 & (\S\ref{sec:i14}) & stands\\\midrule
card \texttt{network-law-free-fermion} & (\S\ref{sec:i16}) & pass\\
card \texttt{corner-kernel-closed-form} & (\S\ref{sec:i16}) & pass\\
card \texttt{network-global-bound-coincident} & range of (1) (\S\S\ref{sec:i9},\ref{sec:i16}) & \textbf{fail}\\
card \texttt{network-upper-uniform-two-corners} & (\S\ref{sec:i16}) & pass\\
card \texttt{network-law-c-linear} & (\S\ref{sec:i16}) & pass\\
card \texttt{network-remainder-uniformity} & (\S\S\ref{sec:i15},\ref{sec:i16}) & pass\\\bottomrule
\end{tabularx}
\end{center}

\section{Item 1: frame reduction and canonical composite (lem:frame, lem:canon)}\label{sec:i1}
\lstep{1}{1}{lem:frame $\langle1\rangle1$--$\langle1\rangle3$ are correct.}
\lby{recomputed: $g=(2x-p_1-p_{n+1})/(x-p_1)$ has $ad-bc=p_{n+1}-p_1>0$, $1-g=(p_{n+1}-x)/(x-p_1)$, so
 $y\circ g=y_{I_1}$; differentiating gives $g'\beta_{I_1}=\beta_0\circ g$; the composition rule
 $\Sch(f\circ h)=\Sch(h)+(h')^2\Sch(f)\circ h$ (brief, ``Setting'') applied to $\Phi\circ g^{-1}$ turns
 $\delta_{p_k}\circ g^{-1}$ into $g'(p_k)\delta_{g(p_k)}$ and multiplies by $(g^{-1})'(g(p_k))^2=g'(p_k)^{-2}$, hence
 the mass $\kappa_k/g'(p_k)$, and $\Sch(g)=0$ adds nothing; the strength is then $\kappa_k\beta_{I_1}(p_k)$}
\lstep{1}{2}{lem:frame $\langle1\rangle4$ and lem:canon are correct.}
\lby{$\beta_{V_{\Phi^g}}\circ\beta_{W_g}=\beta_{W_gV_\Phi}=\beta_{W_g}\circ\beta_{V_\Phi}$ and $\omega\circ\beta_{W_g}=\omega$
 (G1a lem:moebius-inv, $W_g$ unitary on all of $H$ although the pole of $g$ is the end $p_1$); lem:canon
 $\langle1\rangle2$: every factor with index $j>k$ is the identity on $(-\infty,q_j]\ni q_k$, so $G_k'(q_k)=1$;
 $\langle1\rangle4$ is G1a lem:same-schwarz plus Möbius invariance}
\lstep{1}{3}{Hypothesis (H) is stated ((C8), with G1a def:H) and used exactly once: to know the protocol's
 Schwarzian (G1a thm:corner), after which only the corner measure enters.}
\lby{(C8), lem:canon last sentence, thm:T1/thm:T2 last sentences.  The scaled family is \emph{defined} in (C8) as the
 corner-measure family $s\sum\kappa_k\delta_{p_k}$ (the brief's definition), so its state is $\omega_s$ by lem:canon
 and lem:frame without rem:scaled; rem:scaled is an additional interpretation (protocol-shaped composite with
 $\sigma_m\to s\sigma_m$) that nothing below uses.  G1a thm:state(1) (which needs admissible steps, $s_m\ge\lambda_m$)
 is \emph{not} used for $s<1$: normality of $\omega_s$ comes from PA's implementer (\S\ref{sec:i2})}
\lstep{1}{4}{``Right compressions in canonical order'' and ``the protocols'' are the same object at the level of
 \emph{states} for the free fermion, not at the level of maps.}
\lby{$\Psi_s$ and a protocol's composite differ by a post-composed Möbius map (G1a lem:same-schwarz), which does not
 change the quasi-free state (G1a thm:state(2)); for a general net this identification needs the implementer-level
 Möbius step, which \doc\ correctly lists as open (T6 verdict box; \S\ref{sec:i14})}
\begin{verdictok}
lem:frame, lem:canon: STAND.  (H) is stated and used; the identification with protocols is the Schwarzian
dependence of the free-fermion state, exactly as claimed.
\end{verdictok}

\section{Item 2: the implementer (lem:implementer)}\label{sec:i2}
\lstep{1}{1}{(i) holds: $d_k\circ\kt_{s\zeta_k}\circ d_k^{-1}=\Rp_{q_k,s\kappa_k}$ on $I$.}
\lby{$d_k(x)=q_k+(1-q_k)x$ is affine with $d_k(0)=q_k$; conjugating $x\mapsto x/(1+\sigma x)$ ($x>0$) gives
 $x'\mapsto q_k+(x'-q_k)/(1+\sigma(x'-q_k)/(1-q_k))$, i.e.\ parameter $s\zeta_k/(1-q_k)=s\kappa_k$; PA lem:ext(E2)
 gives $\kt_\sigma|_I=\Rp_{0,\sigma}|_I$; order of the product $W^{(1)}\cdots W^{(M)}=V_{\Rp_1\circ\dots\circ\Rp_M}$ on
 $H_I$ matches eq.~(Psi)}
\lstep{1}{2}{(ii) telescoping and (iii) the bound $\norm{G_{\rm tot}(s)}_2\le C_0s\sum\zeta_k$ hold \emph{for every}
 $s\ge0$.}
\lby{the identity $\Pp-W_2^*W_1^*\Pp W_1W_2=(\Pp-W_2^*\Pp W_2)+W_2^*(\Pp-W_1^*\Pp W_1)W_2$ (checked); the
 per-factor bound $\norm{G_\sigma}_2\le C_0|\sigma|$ for all real $\sigma$ is PA's Correction C3 remark
 (l.~1606--1612: $\norm{G_s}_2\le|s|\norm{G^{(1)}}_2\le C_0|s|$ from \texttt{lem:integralrep}), not PA prop:HS, whose
 own range is $s\norm\chi_\infty\le\frac12$; [K3](b) of \doc\ cites both, correctly}
\lstep{1}{3}{(iv): every hypothesis of PA cor:implementer and prop:vector is met.}
\lby{both off-diagonal blocks of $W_{\rm tot}(s)$ are Hilbert--Schmidt by (iii) (PA Correction C4 requires both,
 and both are supplied); $W_{\rm tot}f=V_{\Psi_s}f$ on $H_I$ by (i)}
\lstep{1}{4}{Note~4 Thm.~7.1 and G1a thm:state(1) are not needed and not used for $s<1$.}
\lby{the scaled steps are not admissible for $s<1$ (G1a lem:admissible), which \doc\ says in (C8); normality comes
 from (iv).  At $s=1$ the state is the protocol's (\S\ref{sec:i1})}
\begin{verdictok}
lem:implementer: STANDS.  (v) $\delta(s)=E_IG_{\rm tot}E_I$ follows from (i) as written.
\end{verdictok}

\section{Item 3: the first-order defect of the composite (lem:Dtot)}\label{sec:i3}
\lstep{1}{1}{The frame of the translation is $+y_k$: $E_IW_{d_t}G^{(1)}W_{d_t}^*E_I=D_t$ with $y$-kernel
 $\widetilde D(y-t,y'-t)$, and the corner of $D^{(1)}$ sits at $y=0$.}
\lby{$U_0V_{d_t}U_0^{-1}=\tau_t$, $(\tau_tg)(y)=g(y-t)$ (\doc\ eq.~(transl), recomputed: $(d_t^{-1})'=e^t$,
 $\beta_0(x')=e^{-y'}$); $d_{y_k}(0)=q_k$}
\lstep{1}{2}{The sign of eq.~(ddot) is the one forced by PA eq.~(hardy).}
\lby{with $\Pp(x,y)=\frac12\delta-\frac i{2\pi}\mathrm{pv}\frac1{x-y}$ the recovered symbol has kernel
 $\sqrt{\Psi'(x)\Psi'(x')}\Pp(\Psi(x),\Psi(x'))$, so $\delta(s)=\frac i{2\pi}\kappa_{\Psi_s}$ (PA lem:kernel); for one
 corner at $0$, $\partial_\sigma\kappa|_0=xx'/(x-x')^2$ for $x<0<x'$, and in the $y$-frame (half-density factor
 $e^{-(y+y')/2}$, $x=2e^{-y/2}\sinh\frac y2$, $x-x'=2e^{-(y+y')/2}\sinh\frac{y-y'}2$) this is
 $\frac{\sinh(y/2)\sinh(y'/2)}{\sinh^2((y-y')/2)}$; with the factor $-\frac i{2\pi}$ from $\delta=Q_I-\Qt$ this is
 eq.~(ddot), and it is \texttt{f3\_t0.py}'s \texttt{ddot} (its \texttt{kern} $=-\sinh\sinh/\sinh^2$ times
 $\frac i{2\pi}$ in the $A$--$D$ block)}
\lstep{1}{3}{$\langle1\rangle2$--$\langle1\rangle5$ of lem:Dtot are correct.}
\lby{PA prop:expansion $\langle1\rangle2$ (in its range $s\norm\chi_\infty\le\frac12$, enough for $s\downarrow0$);
 strong convergence $U_k(s)\to\one$ and the finite-rank density argument; $\norm{E_IXE_I}_2\le\norm X_2$}
\lstep{1}{4}{Independent numerical instance (item L of \S\ref{sec:num}): for $M=2$, $(q,\kappa)=((-0.5,0.9),(0.4,1.7))$
 and $M=3$, $(q,\kappa)=((-1.5,0.3),(0.2,1.2),(0.75,2.5))$, at 8 points each, the complex-step derivative of the exact
 kernel equals $\sum_k\zeta_k\widetilde D(y-y_k,y'-y_k)$ to relative $2.0\cdot10^{-50}$ (tolerance $10^{-12}$), while
 the wrong frame $\widetilde D(y+y_k,y'+y_k)$ is off by $133\,\%$.}
\lby{\texttt{rigor/ref\_p6\_4\_composite.out}, block L; my configurations, points and differentiation rule (complex
 step) differ from \texttt{g1b\_kernel.py}'s (Richardson)}
\begin{verdictok}
lem:Dtot: STANDS, with sign and frame $y_k$ (not $-y_k$) confirmed.  lem:polar: STANDS ($\overline{D_a}D_b=
\abs{D^{(1)}}^2e^{-iu(b-a)/2\pi}$ by eq.~(Dt); the bound $\iint\abs{\Dtot}^2/N\le Z^2/(3\pi^2)$ is Minkowski with QL
prop:gB, $2r\iint\abs{D^{(1)}}^2/N=2r/(3\pi^2)$).
\end{verdictok}

\section{Item 4: differentiability of the compressed symbol (lem:polar, lem:taylor)}\label{sec:i4}
\lstep{1}{1}{The exact hypothesis of QL thm:lower that the composite lacks is QL lem:standing(c): $z\mapsto R_z(u,v)$
 holomorphic on $\{|z|<1\}$ with $|R_z|\le|z|/(L(1-|z|))$, used at thm:lower $\langle1\rangle4\langle2\rangle1$ to make
 $\zeta\mapsto Q_P-P\delta_\zeta P$ holomorphic and hence $\Phi_P$ real-analytic ($\langle2\rangle4$), which feeds
 the Taylor step $\langle1\rangle5$.}
\lby{QL l.~135--150 and l.~437--470, read}
\lstep{1}{2}{lem:taylor replaces it correctly: analyticity is moved from the curve to the function of the symbol.}
\lby{$f(Z)=-\log F(\omega_X,\omega_Z)$ is real-analytic on the open set $\{0<Z<1\}$ of a finite-dimensional real
 space (QL thm:lower $\langle2\rangle3$, Note~5 Thm.~2.1: $F=\det[(1-X)(1-Z)]^{1/2}\det[1+(G_X^{1/2}G_ZG_X^{1/2})^{1/2}]>0$);
 $X$ is an interior minimum, so $\mathrm Df(X)=0$; the Bures box is applied only to the straight lines $X+\eps Z$
 (real-analytic, faithful); Taylor's theorem with remainder for a $C^3$ function in finitely many variables gives
 $f(Y_s)=\frac12\mathrm D^2f(X)[H_s,H_s]+O(\norm{H_s}^3)$ with $H_s=s\dot Y+o(s)$.  Only $\norm{Y_s-X-s\dot Y}=o(s)$
 is used, which lem:Dtot supplies in Hilbert--Schmidt norm}
\lstep{1}{3}{The identification $\mathrm D^2f(X)[Z,Z]=\frac14g_B$ and the Legendre inequality are correct.}
\lby{comparison of $\eps^2$-coefficients of the Taylor expansion and of the Bures box ($-\log F=\frac{\eps^2}8g_B+
 O(\eps^3)$, QL eq.~(bures)); QL prop:oneparticle needs only a differentiable curve of symbols in $(0,1)$ with hermitian
 derivative, which the line $X+\eps\dot Y$ is}
\begin{verdictok}
lem:taylor: STANDS.  First-order differentiability of the compressed symbol is enough; no property of $s\mapsto\delta(s)$
beyond $\norm{\delta(s)-s\Dtot}_2=o(s)$ is used.  The spectrum of $Y_s$ stays in $(0,1)$ for every $s$ (T1 $\langle1\rangle4$),
so the QL step $\langle2\rangle2$ is not needed either.
\end{verdictok}

\section{Item 5: the lower bound (thm:T1), direction $\ge$}\label{sec:i5}
Direction ledger: quantity $\liminf_{s\downarrow0}\Phi_{\rm tot}(s)/s^2$ (exact functional, free fermion, $\cF(I)$);
evidence: $\Phi_{\rm tot}\ge\Phi_P$ (lower bound, exact), $\lim\Phi_P/s^2=\frac18g_B^{(P)}$ (equality, finite
dimensions), $g_B^{(P)}\ge\frac18\cdot$ Legendre bracket (lower bound); all three point the same way.
\lstep{1}{1}{Data processing is used in the right direction.}
\lby{restriction $\cF(I)\to\cM_P$ is a channel, fidelity does not decrease, $-\log$ reverses: $\Phi_{\rm tot}\ge\Phi_P$
 (QL thm:lower $\langle1\rangle3$); $\omega_s$ is normal by lem:implementer(iv), so the restriction is defined}
\lstep{1}{2}{The bound is linear in the defect: only $\dot Y=-P\Dtot P$ enters, and $Y_s-X=-P\delta(s)P$.}
\lby{eq.~(Qts): the symbol of $\omega_s$ is $Q_I-\delta(s)$; lem:Dtot}
\lstep{1}{3}{The compression $P$ is legitimate and the sign of the trial operator is handled.}
\lby{$0<X,Y_s<1$ by QL lem:standing(a) for $Q_I$ and for $Q_{\Psi_s(I)}$ ($\Psi_s(I)$ is a half-line $(-\infty,\Psi_s(1))$,
 its complement has positive measure), pulled back by the unitary $V_{\Psi_s}:H_I\to H_{\Psi_s(I)}$; with $\ell'=-\ell$ the
 bracket is $2\Tr(\Dtot\ell)-\Tr(Q_I\ell(1-Q_I)\ell)=\Omega_{\rm tot}(\ell)$; $\sup_\ell\Omega_{\rm tot}=2\iint\abs{\Dtot}^2/N$
 needs only hermiticity of $\Dtot$ ([K2](i), QL prop:gB $\langle1\rangle2$--$\langle1\rangle5$)}
\lstep{1}{4}{$r$ components: fidelities multiply, $g_Q$ carries the factor $r$.}
\lby{PA thm:main-upper $\langle1\rangle1$, as cited}
\begin{verdictok}
thm:T1: STANDS for every finite configuration (C6), with the constant $g_Q(\Dtot)=c\sum\zeta_j\zeta_k\Gh(y_j-y_k)$.
\end{verdictok}

\section{Item 6: the upper bound (lem:prodexp, prop:inf, thm:T2), direction $\le$}\label{sec:i6}
Direction ledger: $\Phi_{\rm tot}\le-\log\abs{\ip\Om{\Gamma(U_s)\Om}}$ (upper bound: Uhlmann with one trial vector),
$\le\norm{V_s^{\rm tot}}_2^2/(2(1-\norm{V_s^{\rm tot}}^2))$ (upper bound, PA cor:master), then the infimum over $h'$
(an infimum of upper bounds is an upper bound).  All in the same direction.
\lstep{1}{1}{lem:prodexp: the expansion $\norm{W_{\rm tot}^*V_s^{\rm tot}-s(\Pm\Dc_{\rm tot}\Pp+iN')}_2=o(s)$ is
 correct, and both blocks of $U_s$ go to $0$ in norm.}
\lby{$W_{\rm tot}^*\Pm W_{\rm tot}=\Pm+G_{\rm tot}$; insertion of $\Pp+\Pm$; PA prop:expansion $\langle1\rangle3$ for
 $Z_s$; lem:Dtot $\langle1\rangle4$ times $\Pp$ with $T_k\Pp=\Pm\Dc_k\Pp$ ($W_{d_k}$ commutes with $P_\pm$); for
 $\Pp U_s\Pm$ the decomposition through $\Pp W_{\rm tot}\Pm=-W_{\rm tot}G_{\rm tot}\Pm$; all rechecked}
\lstep{1}{2}{PA's Uhlmann step and the overlap formula apply to the product.}
\lby{PA prop:uhlmann uses of the unitary only that it restricts to $V_{\Psi_s}$ on $H_I$ and has an implementer
 (both blocks Hilbert--Schmidt, lem:prodexp(iii)); PA thm:det needs $\norm V<1$ \emph{and} $\norm{\Pp U\Pm}<1$, and
 T2 $\langle2\rangle3$ checks both (the omission of the second that RR Correction C3 records for RR is not repeated
 here); with both blocks of norm $<1$, $\Pp U\Pp$ is invertible, so no index question arises}
\lstep{1}{3}{prop:inf: PA prop:infimum $\langle1\rangle1$--$\langle1\rangle6$ and lem:tangentvec use the generator
 only through $\Pm\Dc\Pp\in\Stwo$, $G^{(1)}=[\Pm,\Dc]$ with $\Pp\Dc\Pm=-(\Pm\Dc\Pp)^*$, and the value of
 $\sup_\ell\Omega$.}
\lby{read at PA l.~1096--1160 and 1216--1257: lem:wick and thm:three-faces do not involve $\Dc$; for $\Dc_{\rm tot}$,
 $[\Pm,\Dc_{\rm tot}]=\sum\zeta_kT_k$ and the skew relation hold factor by factor after conjugation by $W_{d_k}$;
 the bookkeeping $\varphi-\mathrm{Var}=\frac14\Omega_{\rm tot}(-2\ell)$ was recomputed:
 $\frac14\Omega_{\rm tot}(-2\ell)=-\Tr(\Dtot\ell)-\Tr(Q\ell(1-Q)\ell)$}
\lstep{1}{4}{Constant check at $M=1$: $\inf_{h'}\norm{\Pm(\Dc+ih')\Pp}_2^2=2g_Q(D^{(1)})=2\cdot\frac14\cdot\frac1{3\pi^2}
 =\frac1{6\pi^2}$, PA eq.~(infimum); hence $\limsup\le\frac1{12\pi^2}$, Theorem~A.}
\lby{prop:inf at $M=1$, lem:polar, QL prop:gB}
\begin{verdictok}
lem:prodexp, prop:inf, thm:T2: STAND.  Every hypothesis of PA thm:main-upper's chain (implementer with both blocks,
vector representative, commutant trial unitaries $e^{ish'}$ of the form $\one\oplus W'$, overlap formula with both
norms $<1$, three faces, Wick reduction) is checked for the multi-corner product; the multi-jump field never needs an
implementer of its own.
\end{verdictok}

\section{Item 7: the law at fixed configuration (cor:law)}\label{sec:i7}
\lstep{1}{1}{cor:law asserts the limit at fixed corner data only, and says so.}
\lby{statement: ``$(s\downarrow0$, fixed corner data$)$'' and ``The $o(s^2)$ is \emph{not} asserted to be uniform'';
 T1 has no rate (finite-rank $P$ chosen after $\ell$), T2's $o(s)$ comes from strong continuity}
\lstep{1}{2}{No uniformity is claimed by accident.}
\lby{every occurrence of ``uniform'' in \doc\ (\texttt{grep -n uniform}: l.~32, 637, 912, 927--1054, 1206--1223)
 either denies uniformity or refers to thm:T5 (upper half, $M=2$) and its use in cor:vwz; cor:W3W4(ii) says ``again
 at fixed configuration''; the Summary table marks T5 ``partial''}
\begin{verdictok}
cor:law: STANDS, at fixed configuration, as stated.
\end{verdictok}

\section{Item 8: no one-parameter group (lem:nogroup)}\label{sec:i8}
\lstep{1}{1}{$\langle1\rangle1$--$\langle1\rangle3$ are correct.}
\lby{$\Psi_t(q_2)=\Rp_{q_1,t\kappa_1}(q_2)\in(q_1,q_2)$; $\Psi_s$ is Möbius between its breakpoints, so the jump of
 $(\Psi_s\circ\Psi_t)''/(\Psi_s\circ\Psi_t)'=(\Psi_s''/\Psi_s')\circ\Psi_t\cdot\Psi_t'+\Psi_t''/\Psi_t'$ at $q_2$ is that of
 $\Psi_t''/\Psi_t'$, i.e.\ $-2t\kappa_2$; push-forwards of all $f\in H_I$ determine the map}
\lstep{1}{2}{Exact recomputation (item N of \S\ref{sec:num}), $q=(-\frac12,\frac13)$, $\kappa=(\frac32,\frac54)$,
 $s=\frac27$, $t=\frac35$, one-sided jets in rational arithmetic: mass of $\Sch(\Psi_s\circ\Psi_t)$ at $q_2$ is
 $-\frac32=-2t\kappa_2$, that of $\Sch(\Psi_{s+t})$ is $-\frac{31}{14}=-2(s+t)\kappa_2$; at $q_1$ both are
 $-\frac{93}{35}=-2(s+t)\kappa_1$; here $\Psi_t(1)=\frac{73}{774}<q_2$, so the corner of $\Rp_{q_2,s\kappa_2}$ lies
 outside $\Psi_t(I)$ and contributes nothing (G1a thm:general).}
\lby{\texttt{rigor/ref\_p6\_4\_composite.out}, block N}
\begin{verdictok}
lem:nogroup: STANDS.  The Schwarzian mass $-2t\kappa_2$ is exact.
\end{verdictok}

\section{Item 9: the global bound (thm:T3), direction $\le$}\label{sec:i9}
\lstep{1}{1}{$\langle1\rangle2$--$\langle1\rangle4$ and $\langle1\rangle6$--$\langle1\rangle7$ are correct.}
\lby{RR prop:rrr $\langle1\rangle1$--$\langle1\rangle2$ hold for every $s\ge0$ (Dyson cocycle; both blocks via PA C2);
 moving $c^{(k)}$ to the right: $X^*c^{(k)}X$ fixes $H_I$ because $XH_I\subseteq H_I$; $\Pm AB\Pp=\Pm A\Pp B\Pp+
 \Pm A\Pm B\Pp$; $\det(\one-V^*V)\ge1-\norm V_2^2$; the infimum $2\hat f_2$ is PA eq.~(infimum)}
\lstep{1}{2}{\textbf{GAP} at $\langle1\rangle5$.  Objected sentence: ``RR Lemma~\texttt{lem:groupbound} for
 $e^{\pm s\zeta_kK}$ \dots: $\norm{\Pm e^{\sigma K}\Pp}_2\le\abs\sigma\norm{\Pm K\Pp}_2$ \emph{for all real} $\sigma$''
 (and the verbatim box [K4], ``for all $s$'').}
\lby{RR's Correction C3 (l.~1026--1079), part (ii) of ``Status of the uses'': the equality $\norm{G_s}_2^2=2\norm{V_s}_2^2$
 behind this bound needs $\operatorname{ind}(\Pp e^{sK}\Pp)=0$, which ``the documents settle \dots\ only while
 $\norm{G_s}_2<1$'', i.e.\ for $\abs s\norm{[\Pm,K]}_2<1$; ``in Lemma~6.1 [\texttt{lem:groupbound}], which is stated for all
 $s$, every $s$ with $\abs s\norm{[\Pm,K]}_2\ge1$'' is open (card \texttt{group-index-zero-large-s}, conjectural).
 Without index zero only $\norm{\Pm e^{\sigma K}\Pp}_2\le\norm{[\Pm,e^{\sigma K}]}_2\le\abs\sigma\norm{[\Pm,K]}_2=
 \sqrt2\abs\sigma\norm{\Pm K\Pp}_2$ is available, a factor $\sqrt2$ too weak}
\lstep{1}{3}{Size of the gap.  For $h'$ near the infimum, $\norm{[\Pm,K]}_2^2=2\norm{\Pm K\Pp}_2^2\downarrow4\hat f_2$,
 so the settled range for factor $k$ is $s\zeta_k<(4\hat f_2)^{-1/2}=\pi\sqrt3\approx5.44$, while T3 claims every
 $s$ with $sZ<(2\hat f_2)^{-1/2}=\pi\sqrt6\approx7.70$.}
\lby{$\hat f_2=1/(12\pi^2)$.  For $M=1$ (allowed by (C6)) the uncovered range is RR's own open range
 $\pi\sqrt3\le\zeta<\pi\sqrt6$, recorded on card \texttt{rate-of-remainder}; for $M\ge2$ it is empty iff
 $\max_k\zeta_k\le Z/\sqrt2$, and nonempty e.g.\ for $\zeta=(1,0.01)$: $s\in[5.44,7.62)$}
\lstep{1}{4}{Repair (either suffices).  (a) Add the hypothesis $s\max_k\zeta_k<\pi\sqrt3$ (equivalently
 $4\hat f_2s^2\zeta_k^2<1$ for every $k$): choose $h'$ with $\norm{\Pm K\Pp}_2^2<\min\{(sZ)^{-2},(2s^2\max\zeta_k^2)^{-1}\}$
 close to $2\hat f_2$; then $s\zeta_k\norm{[\Pm,K]}_2<1$ for every factor, index zero holds by RR C3(ii), and
 $\langle1\rangle5$--$\langle1\rangle7$ go through verbatim.  (b) Keep the range and make the excess conditional on
 \texttt{group-index-zero-large-s} (theorem statement and card \texttt{depends}).}
\lby{RR C3(ii); $\langle1\rangle2$, $\langle1\rangle3$}
\lstep{1}{5}{The comparison in the verdict box is correct, in both directions.}
\lby{$0<\Gh(\Delta)<\Gh(0)$ for $\Delta\ne0$ (T4(d),(e); my grid check, \S\ref{sec:num} B, E) and $\zeta_k>0$ give
 $\sum\zeta_j\zeta_k\Gh(y_j-y_k)\le\Gh(0)Z^2$, equality iff all $y_j$ coincide; $x\mapsto-\frac12\log(1-2x)$ is
 increasing, so the conjectured network form would be the \emph{smaller} (stronger) bound, T3 the weaker; and
 ``implies T2's upper half only'' is the right relation (expansion of $-\frac12\log(1-2x)$)}
\begin{verdictgap}
thm:T3: STANDS WITH REPAIR (repair (a) or (b) of $\langle1\rangle4$).  Not a BREAK: nothing shows the bound false in
the uncovered range.  The verdict box's ``not proved'' for the network-form bound and the relation ``implies only''
are accurate.
\end{verdictgap}

\section{Item 10: the kernel $\Gh$ in closed form (thm:T4)}\label{sec:i10}
\paragraph{(a), (b).}
\lstep{1}{1}{(a) $\langle1\rangle1$--$\langle1\rangle3$: the constants are right.}
\lby{recomputed: $\dd\kappa\dd\kappa'=\frac12\dd u\dd v$ gives $\frac r8\int\cos\frac{u\Delta}{2\pi}S(u)\dd u$ with
 $S$ of QL lem:V ($\frac23\tanh\frac u2/(u(u^2+4\pi^2))$); $u=2\pi q$: $\frac18\cdot\frac23\cdot2\pi\tanh(\pi q)/(8\pi^3q(1+q^2))
 =\tanh(\pi q)/(48\pi^2q(1+q^2))$; in $p$: $\abs{\hat d_0}^2=4\pi^2\abs{D^{(1)}}^2$, $\abs{\dd v/\dd p}=4\pi$, so
 $\frac14\cdot4\pi^2\cdot\frac1{4\pi}S(2\pi q)=\rho(q)$}
\lstep{1}{2}{Independent check of $\rho(q)=\frac14\int\abs{\hat d_0(p,p-q)}^2/W(p,p-q)\dd p$ at $q=0.1,0.5,1,2,5$:
 maximal relative deviation $2.7\cdot10^{-51}$ (tolerance $10^{-10}$).}
\lby{item R of \S\ref{sec:num} (mpmath tanh-sinh, 160 digits, $\hat d_0$ and $W$ literally from T4(c) and (C4),
 $\abs p\le45+\abs q$, neglected weight $<10^{-115}$)}
\lstep{1}{3}{(a) $\langle1\rangle4$--$\langle1\rangle6$ and (b) are correct: $\int\rho=\frac1{12\pi^2}$, and strict
 positive-definiteness follows from $\rho>0$ and the identity $\sum c_j\bar c_k\cos q(\Delta_j-\Delta_k)=\frac12(\abs{\sum
 c_je^{iq\Delta_j}}^2+\abs{\sum c_je^{-iq\Delta_j}}^2)$.}
\lby{identity expanded; QL lem:U $\langle1\rangle6$; my $\int_0^\infty\tanh(\pi q)/(q(1+q^2))\dd q=2$ to $8.9\cdot10^{-16}$ at
 three cut-offs $Q=20,40,80$ (item V2)}
\paragraph{(c) and the Fourier convention.}
\lstep{1}{4}{Both forms of $\hat d_0$ follow from KI thm:main by the substitution of (C3).}
\lby{$\kappa=-2\pi p$: $q_\kappa=w(p)$, $u=-2\pi(p-p')$, $v=-2\pi(p+p')$, $u^2+4\pi^2=4\pi^2(1+(p-p')^2)$,
 $\sinh\frac u2=-\sinh\pi(p-p')$; both members recomputed}
\lstep{1}{5}{The convention is (C3): $\hat A(p,p')=(2\pi)^{-1}\iint e^{ipy}\widetilde A(y,y')e^{-ip'y'}$, $p=-\kappa/2\pi$.
 It is the one in which \texttt{f3\_t0.py}'s $Q$ is the multiplier $w(p)=(1+e^{-2\pi p})^{-1}$ of its docstring.}
\lby{\texttt{f3\_t0.py} never Fourier-transforms: it uses \texttt{kkt\_continuum.Qmat}, whose kernel is
 $\frac12\delta-\frac i{4\pi}\mathrm{pv}\,\sinh^{-1}\frac{y-y'}2$ (its \texttt{F2} has $F_2'(t)=2\log\tanh\frac{\abs t}4$,
 $F_2''=1/\sinh\frac t2$, $F_2(\pm\infty)=\mp\pi^2$; checked by differentiating the dilogarithm form), and its
 \texttt{ddot} is eq.~(ddot) (\S\ref{sec:i3}); with $\int_0^\infty\sin(pz)/\sinh\frac z2\,\dd z=\pi\tanh\pi p$ (item V3)
 the $e^{+ipy}$ transform of that kernel is $\frac12+\frac12\tanh\pi p=w(p)$.  Under the opposite sign convention
 $\hat d_0$ would change sign ($\hat d_0(-p,-p')=-\hat d_0(p,p')$)}
\lstep{1}{6}{Direct 2-D Fourier integral of eq.~(ddot), my code, at
 $(p,p')=(0.4,-0.1),(1.2,0.3),(-0.9,0.6),(0.35,0.35),(2,-1.5),(0.05,0.8)$: maximal $\abs{\text{direct}-\text{T4(c)}}=
 7.2\cdot10^{-18}$ (tolerance $10^{-9}$), imaginary parts $0$; the sign-flipped formula misses by $\ge3.4\cdot10^{-3}$.}
\lby{item K of \S\ref{sec:num}: $(t,u)$ coordinates on both quadrants, composite Gauss--Legendre with 16 and 24 nodes
 per panel (difference $\le5.8\cdot10^{-17}$), $t\le100$ (integrand $\le100e^{-50}$ beyond); the diagonal point tests the
 removable singularity, where T4(c) is $w'(p)p/\pi$}
\begin{verdictok}
T4(a), (b), (c): STAND; the Fourier convention is (C3), used consistently, and agrees with \texttt{f3\_t0.py}.
\end{verdictok}
\paragraph{(cf), (d), (e).}
\lstep{1}{7}{Every step of the derivation of (cf) is correct.}
\lby{re-derived: Mittag-Leffler (QL lem:U $\langle1\rangle3$, refereed, Ahlfors locator there) divided by $q(1+q^2)$;
 Tonelli for the nonnegative series ($\sum_n\int=\frac\pi2\cdot4$); partial fractions
 $\frac1{(q^2+a^2)(q^2+1)}=\frac1{1-a^2}(\frac1{q^2+a^2}-\frac1{q^2+1})$; lem:lorentz (its Gaussian regularisation:
 $(\pi/\eps)^{1/2}=2\pi(4\pi\eps)^{-1/2}$, error $a\,\mathbb E\abs z=2a(\eps/\pi)^{1/2}$, checked; numerically item V1);
 $\frac1{1-a_n^2}=-\frac4{(2n-1)(2n+3)}$ telescopes to $0$; the three power series and $\frac1{a(1-a^2)}=\frac1a+
 \frac1{2(1-a)}-\frac1{2(1+a)}$; $2-x^2-x^{-2}=-(x^{-1}-x)^2$; $\artanh x=\frac12\log\coth\frac{\abs\Delta}4$ for
 $x=e^{-\abs\Delta/2}$}
\lstep{1}{8}{Numerically, closed form, series and an independent quadrature of $\int\rho(q)\cos(q\Delta)\dd q$ agree at
 $\Delta=0.5,1,1.5,2,3,4,5,6,8$: $\max\abs{\text{QG}-\text{CF}}/\Gh(0)=7.2\cdot10^{-16}$ (tolerance $10^{-8}$), stable under
 $Q=20,40,80$; $\max\abs{\text{SR}-\text{CF}}/\Gh(0)=5.0\cdot10^{-49}$ (tolerance $10^{-12}$).  The values $\Gh/\Gh(0)$ of
 \doc's table and of the card agree with mine in every printed digit (e.g.\ $0.898386580099973$, $0.024419213316489$).}
\lby{item P of \S\ref{sec:num}; QG = composite Gauss--Legendre on $[0,Q]$ plus the \emph{exact} tail by $\mathrm{Ci}$ and
 $E_1$ (not \texttt{quadosc}, not QUADPACK), SR = term-by-term sum, CF at 50 digits}
\lstep{1}{9}{(d): both bounds and the sharpness hold.}
\lby{proof re-derived ($b_n>0$, $\sum_{n\ge1}b_n=\frac23$ from (b), $x^{2n+1}\le x^3$); grid
 $\Delta=0.01,\dots,20$ (2000 points, 50 digits): 0 violations of $x(4-x^2)/(36\pi^2)\le\Gh\le x/(9\pi^2)$ and of
 $\abs{\Gh-x/(9\pi^2)}\le x^3/(36\pi^2)$, smallest relative margin $1.0\cdot10^{-10}$ (at $\Delta=20$, where the bounds
 are asymptotically tight); $9\pi^2e^{\Delta/2}\Gh-1=-9.1\cdot10^{-6},-4.1\cdot10^{-10},-8.5\cdot10^{-19}$ at
 $\Delta=10,20,40$, so the rate $\frac12$ and the constant $\frac1{9\pi^2}=\frac43\Gh(0)$ are sharp (item B)}
\lstep{1}{10}{(e): positivity and strict decrease are proved and confirmed (0 sign changes, 0 non-decreasing steps on the
 grid, item E).  ``$C^1$ but not $C^2$ at $0$'' is true but stated only in Remark~\texttt{rem:kernelshape}, not in
 T4(e), with a prose argument.}
\lby{$C^1$, $\Gh'(0)=0$: M-test, as written.  Not $C^2$: if $\Gh$ were $C^2$ near $0$, Taylor and $\Gh'(0)=0$ would give
 $(\Gh(\Delta)-\Gh(0))/\Delta^2\to\frac12\Gh''(0)$, but by the closed form this ratio is
 $\frac1{12\pi^2}[\frac18-\frac14\log\frac4{\abs\Delta}+o(\log\frac1{\abs\Delta})]\to-\infty$; equivalently, termwise,
 $\Gh''(\Delta)=\frac1{24\pi^2}\sum_na_ne^{-a_n\Delta}/(1-a_n^2)\le\frac1{24\pi^2}[\frac23e^{-\Delta/2}-\sum_{n\ge1}e^{-a_n\Delta}/a_n]
 \to-\infty$.  Numerically $\Gh''(\Delta)-\Gh''(10\Delta)\to-\frac{\ln10}{24\pi^2}=-9.720860\cdot10^{-3}$ (7 digits at
 $\Delta=10^{-5},10^{-6}$; item E): a logarithmic singularity of $\Gh''$, as the remark says}
\begin{verdictok}
T4 (cf), (d), (e) and lem:lorentz: STAND.  The closed form is derived, not fitted.  MINOR: move ``$C^1$, not $C^2$ at $0$''
from rem:kernelshape into T4(e) with the two-line argument of $\langle1\rangle10$ (the card states it as part of (e)).
\end{verdictok}

\section{Item 11: consequences for W3, W4 (cor:W3W4)}\label{sec:i11}
\lstep{1}{1}{(i) is T4(d) at $x=\sqrt\eta$ divided by $\Gh(0)$: $\sqrt\eta(4-\eta)/3\le\Gh(\log\frac1\eta)/\Gh(0)\le\frac43\sqrt\eta$.}
\lby{$x(4-x^2)/(36\pi^2)\cdot12\pi^2=x(4-x^2)/3$, $x/(9\pi^2)\cdot12\pi^2=\frac43x$}
\lstep{1}{2}{(ii) is conditional on $e^\Delta\ge2/\eps$ and bounds each cross term of the second-order coefficient.}
\lby{G1a thm:sep (l.~1606) gives $e^{\Delta_k}\ge2/\eps$ for adjacent corners of an all-weak L$\to$R ordinary-Petz chain;
 non-adjacent separations are sums of adjacent ones; $\zeta_j\zeta_k\Gh\le\eps^2\cdot\frac1{9\pi^2}(\eps/2)^{1/2}$}
\lstep{1}{3}{What is inferred: W4's hypothesis ``$\Gh(\Delta)\to0$'' and W3's hypothesis ``$\Gh\sim e^{-\Delta/2}$'' are
 proved; the other hypothesis of both (uniform remainder) is explicitly left to \S7.  Nothing beyond T4 is used.}
\lby{statement of cor:W3W4 and the ``Effect on the working results'' paragraph}
\begin{verdictok}
cor:W3W4: STANDS; it infers nothing beyond T4(d) and G1a thm:sep, and keeps ``at fixed configuration''.
\end{verdictok}

\section{Item 12: uniform upper bound for two corners (lem:onegroup, thm:T5), direction $\le$}\label{sec:i12}
\lstep{1}{1}{lem:onegroup (ii), (iii) hold for every $\sigma\ge0$; (i) holds only for $\sigma\norm{[\Pm,K]}_2<1$.}
\lby{(ii): $\Pm U_\sigma\Pp=U_\sigma G_\sigma\Pp$, $G_\sigma=\int_0^\sigma U_\tau^*[\Pm,K]U_\tau\dd\tau$ (sign rechecked:
 $\partial_\tau U_\tau^*\Pp U_\tau=U_\tau^*[\Pp,K]U_\tau$), $[\Pm,K]\Pp=k$, no index needed; (iii): strong continuity.
 (i) rests on RR lem:groupbound, whose range is restricted by RR Correction C3 exactly as in \S\ref{sec:i9}.
 MINOR: T5 uses (i) only as $\sigma=s\zeta_j\downarrow0$}
\lstep{1}{2}{The cross term is controlled uniformly in $y_2-y_1$, including $y_2-y_1\to\infty$.}
\lby{$V=V_1A_2+B_1V_2$; in $\Tr(A_2^*k_{y_1}^*B_1k_{y_2})-\Tr(k_{y_1}^*k_{y_2})=\Tr(k_{y_2}(A_2-\Pp)^*k_{y_1}^*B_1)+
 \Tr(k_{y_1}^*(B_1-\Pm)k_{y_2})$ (recomputed) the diagonal block $A_2$ of corner 2 multiplies only $k_{y_2}$ and $B_1$ of
 corner 1 only $k_{y_1}^*$; $\norm{k_{y_2}(U^{(2)*}-\one)}_2=\norm{k(U^*_{s\zeta_2}-\one)}_2$ because $k_{y_2}$ and $U^{(2)}$
 are conjugated by the \emph{same} $W_{d_{y_2}}$.  The remainders $R_j$ are translation invariant likewise.  The
 trial generator $h'$ is fixed by $\delta$ alone ($\langle1\rangle3$), not by the configuration, and $\tau_t$ commutes
 with the projection $\Pi$ because $W_{d_t}\mathcal A'W_{d_t}^*=\mathcal A'$ ($d_t(I')=I'$).  So no operator of one corner
 ever acts on the first-order vector of the other, and the ``far corner near the moving endpoint'' never meets $k$}
\lstep{1}{3}{$M=2$ is used in $\langle2\rangle1$ (two factors) and in $\langle2\rangle3$ (adjacency of each corner's
 diagonal block to its own $k$); for $M\ge3$ an intermediate corner's block sits between $k_{y_j}^*$ and $k_{y_k}$, which is
 (U2) of the verdict box.}
\lby{expansion of $\Pm U^{(1)}U^{(2)}U^{(3)}\Pp$ written out}
\lstep{1}{4}{MINOR: the side condition of $\langle1\rangle4$, $s(\zeta_1+\zeta_2)\norm k_2<1$, must be supplemented by
 $\sqrt2\,s\max_j\zeta_j\norm k_2<1$ (the index-zero range of the factor bounds used through T3 $\langle1\rangle5$).}
\lby{\S\ref{sec:i9}; with $h'$ fixed and $\mathcal Z$ compact this holds for all $s\le s_\eps$ after shrinking $s_\eps$,
 uniformly in the configuration, so the statement of T5 is unaffected}
\begin{verdictok}
thm:T5: STANDS (uniform in all $y_1<y_2$ and locally uniform in $\zeta$, upper bound only, $M=2$).  lem:onegroup: (i)
should read ``for $0\le\sigma<\norm{[\Pm,K]}_2^{-1}$''.
\end{verdictok}

\section{Item 13: the VWZ ratio (cor:vwz)}\label{sec:i13}
\lstep{1}{1}{The corner data are right, recomputed exactly for $(a,b)=(1,1),(1,3),(2,7),(1,100)$ from G1a prop:1B1A and
 prop:prot2 ($c=b$, $d=a$, $\theta_t=1$): $\zeta_1=\zeta_3=a/b$, $\zeta_2=\frac{a(a+2b)}{2b(a+b)}$, $e^{\Delta_{23}}=
 \frac{a+2b}a=1/\eta$, $\zeta_3^{(J)}=\frac{2a}{a+2b}$, $\hat\zeta_2=\frac{s+2}{2(s+1)}$; the brief's L$\to$R formula
 $\zeta_k=2r_k(1+v_k)/((1+r_k)(e^{\Delta_k}-1))$ gives the same $\zeta_2$.}
\lby{item W of \S\ref{sec:num} (rational arithmetic)}
\lstep{1}{2}{Upper half: T5 with $\mathcal Z=[\frac34,1]\times\{1\}$ applies although $y_3-y_2=\log\frac{s+2}s$ moves with
 $s$, because T5 is uniform in the positions; then $\Gh(\log\frac1\eta)\to0$ (cor:W3W4(i)) and $\hat\zeta_2\to1$ give
 $\limsup\le2$.}
\lby{Protocol 2 is $\Rp_{p_2}\circ\Rp_{p_3}$, canonical order, (H) by G1a prop:prot2; lem:frame carries its data to
 the T0 frame}
\lstep{1}{3}{Lower half: data processing is used in the right direction, and the restricted state is a one-corner state.}
\lby{restriction $\cF(I)\to\cF(J)$, $J=(p_2,p_5)$, can only raise the fidelity, so $\Phi^{(2)}\ge\Phi_A(\zeta_3^{(J)})$;
 on $J$, $\Rp_{p_2,\sigma_2}$ is Möbius with pole $p_2-1/\sigma_2\notin(p_2,\infty)$, so G1a thm:state(2) removes it and
 G1a cor:phiA gives $e^{-\Phi_A(\zeta)}$ with $\zeta=\sigma_3\beta_J(p_3)$; $\Phi_A(\frac{2s}{s+2})/\Phi_A(s)\to1$ by
 Theorem~A}
\begin{verdictok}
cor:vwz: STANDS: $1\le\liminf\Phi^{(2)}/\Phi^{(1)}\le\limsup\Phi^{(2)}/\Phi^{(1)}\le2$ as $a/b\to0$, free fermion,
ordinary Petz.  The limit $2$ is correctly reduced to a uniform lower bound (open).
\end{verdictok}

\section{Item 14: $c$-linearity (thm:T6)}\label{sec:i14}
\lstep{1}{1}{$\langle1\rangle1$: (H1), (H2), (H2$'$) hold for each conjugated factor, with $a_k=\Delta_I^{i\tau_k}a\in\mathcal H_T$.}
\lby{C11 thm:main (l.~792) for TB's $g_{\rm tot}=-\chi\in\mathcal E$ ($L=1$); TB (BW) on the line (l.~88--89):
 $\Delta^{i\tau}$ is $u\mapsto e^{2\pi\tau}u$, $u=1-x$, i.e.\ $d_t$ with $t=-2\pi\tau$, so $\tau_k=-y_k/2\pi$ puts the corner at
 $y_k$ (sign checked against TB's audited convention); $\Delta^{i\tau}\Om=\Om$; TB lem:reduce(3) for $\mathcal H_T$}
\lstep{1}{2}{$\langle1\rangle2$: the product implements the extension of $\Psi_s$ and
 $\norm{U(s)^*\Om-\Om+is\,a_{\rm tot}}=o(s)$.}
\lby{covariance composes; induction over the factors in the order $U(s)^*=U^{(M)*}\cdots U^{(1)*}$, each step using
 $U^{(m+1)}(s)^*\to\one$ strongly, so $s(U^{(m+1)*}a_k-a_k)=o(s)$}
\lstep{1}{3}{$\langle1\rangle3$: TB prop:H3, rem:H2H3, prop:H2 need no group property.}
\lby{read at TB l.~290--410: prop:H3 uses (H1) only to make $\omega_s$ a vector state, (H3) on $K,K^2$, and
 thm:three; rem:H2H3 derives (H3) from the vector expansion alone; prop:H2 is hypothesis (V) of thm:uhllower for
 $\Psi_s=U(s)^*\Om$; (RS) for $\cA(I)$}
\lstep{1}{4}{$\langle1\rangle4$: TB prop:reduce and thm:linear apply to $a_{\rm tot}$.}
\lby{their proofs use $b\in\mathcal H_T$ (closest point $E'b\in\mathcal H_T$) and $\operatorname{Im}\ip\Om b=0$
 ($\ip\Om{\Delta^{i\tau}a}=\ip\Om a=0$); $V^{-1}\Delta^{it}V$ is the push-forward on $\mathcal K$ (thm:linear
 $\langle1\rangle2$), whence $V^{-1}a_{\rm tot}=(c/48\pi^2)^{1/2}\sum\zeta_k[\delta_{\tau_k*}g]$; $\frac c{48\pi^2}\cdot\frac12=
 \frac c{96\pi^2}$; $q$ involves only $(\mathcal K,\Delta_0,J_0,g,\zeta,\tau)$}
\lstep{1}{5}{$\langle1\rangle5$: the free-fermion evaluation gives $q=2\sum\zeta_j\zeta_k\Gh(y_j-y_k)$.}
\lby{TB thm:value $\langle2\rangle2$: $T(f)\Om=\Pm(-i\mathfrak L_f)\Pp$, and $\mathfrak L_{-\chi}=-\Dc$, so $a=i\Pm\Dc\Pp$
 and $a_{\rm tot}=i\Pm\Dc_{\rm tot}\Pp$; TB rem:klein (graded net: thm:linear applies verbatim); prop:inf with
 $\dist(-b,R)=\dist(b,R)$ for the real subspace $R=\overline{\cF(I)'_{\rm sa}\Om}$; $\frac18\cdot4c\,q=c\sum\zeta_j\zeta_k\Gh$}
\lstep{1}{6}{Scope.  The theorem concerns the canonical product eq.~(Ugen) and says so (its $\omega^{\cA}_s$ is defined
 by eq.~(Ugen)); its identification with protocols is restricted to right compressions in the canonical order, where
 the product of step implementers \emph{is} eq.~(Ugen); for other orders it needs $U_{h\circ\Psi}\in\mathbb T\,U(h)U_\Psi
 \cA(I)'$, which C11 \S8(b) (l.~824--829) gives only for integer-$c$ Fock representations.}
\lby{T6 verdict box, C11 l.~824--829 read}
\begin{verdictok}
thm:T6: STANDS for the canonical product under TB's standing hypotheses (S2), (HD), (BW), (RS).  Status word: the
proof is complete and uses only refereed cards (\texttt{thm-b-universality}, \texttt{implementer-hypotheses}, both
refereed) and two proved ones of this document; ``proved'' is the weakest word the evidence supports and also the
strongest it allows before a referee pass; no ``under'' qualifier is needed beyond TB's standing hypotheses, which the
Statement names.  The fixed-configuration restriction is in the Statement.
\end{verdictok}

\section{Item 15: the open card network-remainder-uniformity}\label{sec:i15}
\lstep{1}{1}{(U1), (U2) on the card are the two items of \doc's verdict box closing \S7 (l.~1046--1059), in substance and
 in their quantifiers.}
\lby{card Statement vs.\ \doc\ l.~1050--1058: (U1) uniform lower bound, $M\ge2$, needs a uniform spectral gap of $PQ_IP$
 and a uniform bound on $\norm{P(\delta(s)-s\Dtot)P}$ for translated trial operators; (U2) uniform upper bound, $M\ge3$,
 needs $\sup_{t>0}\norm{k(W_{d_t}U_\sigma W_{d_t}^*-\one)}_2\to0$}
\lstep{1}{2}{They are exact in the sense that they name the step at which each proof route stops (T1's finite-rank
 compression without a rate, lem:taylor's constant $C$ depending on the spectrum of $X$; T5 $\langle2\rangle3$ for
 $M\ge3$), not that no other route exists; the card does not claim more.}
\lby{\S\S\ref{sec:i4}, \ref{sec:i12}}
\lstep{1}{3}{\texttt{next} names a lemma ((U1) for $M=2$ as a new lemma in \S7 of \doc) or a named computation with an
 existing script (\texttt{numerics/networks/corner\_kernel.py}, whose block K4 computes $-\log F$ for two corners at finite
 $\zeta$, l.~507--610).}
\lby{file opened}
\begin{verdictok}
network-remainder-uniformity: status \texttt{open} is right; question well posed (quantifiers: $\lim_{s\to0}\sup$ over
separations and over compact $\zeta$-sets); evidence \texttt{doc:...\#sec:T5} resolves (l.~927).
\end{verdictok}

\section{Item 16: the cards, the findings entries, the notes}\label{sec:i16}
Every evidence locator resolves (\texttt{grep -n}): thm:T1 l.~487, thm:T2 l.~599, cor:law l.~629, lem:nogroup l.~660,
thm:T3 l.~684, thm:T4 l.~737, sec:T5 l.~927, lem:onegroup l.~933, thm:T5 l.~958, cor:vwz l.~1008, thm:T6 l.~1075;
\texttt{rigor/g1b\_kernel.py} exists and reproduces its output (\S\ref{sec:num}).  Every \texttt{depends} card exists
(\texttt{corner-calculus}, \texttt{const-gb}, \texttt{const-f2}, \texttt{bures-three-faces}: proved;
\texttt{thm-a-quadratic-law}, \texttt{thm-b-universality}, \texttt{implementer-hypotheses}: refereed).
\begin{enumerate}[leftmargin=1.5em,label=(\roman*)]
\item \texttt{network-law-free-fermion} [proved]: Statement = cor:law with T1, T2 (constant $c\sum\zeta_j\zeta_k\Gh$,
 ``FIXED configuration'', non-uniformity stated, (H) stated).  For a bounded $I$ the Statement is read through
 lem:frame, which the evidence includes implicitly via cor:law; acceptable.  \texttt{depends} covers G1a, QL/PA, KI,
 three faces.  How to verify reproduces (item L: $3.7\cdot10^{-18}$).  \textbf{Pass.}
\item \texttt{corner-kernel-closed-form} [proved]: (a)--(e) and the five values agree with thm:T4 and with my
 recomputation in every printed digit.  MINOR: ``$C^1$ but not $C^2$ at $0$'' is in rem:kernelshape (l.~898--906), not in
 thm:T4(e); true (\S\ref{sec:i10} $\langle1\rangle10$).  \texttt{depends} (\texttt{const-gb}: KI thm:main, QL prop:gB,
 lem:V; \texttt{const-f2}) complete.  G3's note of 18:13 (box and T0 frames minus \texttt{g1b\_kernel.out}) compares
 against the values I reproduce, so it is consistent with this card.  \textbf{Pass}, with a note.
\item \texttt{network-global-bound-coincident} [proved]: part (1) states thm:T3 for \emph{every} $s$ with
 $2s^2Z^2/(12\pi^2)<1$; proved only where also $s\max_k\zeta_k<\pi\sqrt3$ (\S\ref{sec:i9}).  Part (2) (lem:nogroup)
 stands.  Missing edge: the comparison ``weaker \dots\ because $0<\Gh(\Delta)<\Gh(0)$'' uses T4(d),(e), i.e.\
 \texttt{corner-kernel-closed-form}.  \textbf{Fail} (stands with repair: restrict the range in the Statement,
 or add \texttt{group-index-zero-large-s} to \texttt{depends} and mark the excess conditional; add the edge).
\item \texttt{network-upper-uniform-two-corners} [proved]: Statement = thm:T5 + cor:vwz, quantifiers exact (``for every
 compact $\mathcal Z$ and $\eps$ there is $s_\eps$ \dots\ ALL positions''), open parts stated.  Its edge to
 \texttt{network-global-bound-coincident} uses T3 $\langle1\rangle3$--$\langle1\rangle6$ at small $s$ only, which the range
 gap does not touch.  MINOR notes of \S\ref{sec:i12}.  \textbf{Pass.}
\item \texttt{network-law-c-linear} [proved]: Statement = thm:T6 (standing hypotheses, canonical product, fixed
 configuration, scope and the open Möbius step); status word right (\S\ref{sec:i14}).  \textbf{Pass.}
\item \texttt{network-remainder-uniformity} [open]: \S\ref{sec:i15}.  \textbf{Pass.}
\end{enumerate}
\paragraph{The six entries G1b appended to \texttt{rigor/findings\_entries.tex}} (read with \texttt{git diff}, not
edited): plain \LaTeX\ with the project's \verb|\finding| macro; theorem numbers (Thms.~3.2, 4.3, Cor.~4.4, Thm.~5.2,
6.1, 7.2, Cor.~7.3, Thm.~8.1) match \texttt{network\_second\_order.aux}.  Accurate, except the third (GAP) entry: it
states the proved T3 bound with no range at all; it should carry ``for $2s^2(\sum\zeta_k)^2/(12\pi^2)<1$ and
$s\max_k\zeta_k<\pi\sqrt3$''.
\paragraph{G1b's notes on the orchestrator cards.}  \texttt{corner-correlation-kernel}: accurate in every number and
claim.  \texttt{network-second-order-law}: accurate, except that ``proved instead the coincident-corner bound thm:T3''
needs the range of \S\ref{sec:i9}.  \texttt{vwz-protocol-ordering}: accurate, except one ambiguity: ``T6 makes the
Protocol 2 statement $c$-linear for every net'' holds for the fixed-configuration second-order law of Protocol 2's corner
data only; the asymptotic bracket $1\le\Phi^{(2)}/\Phi^{(1)}\le2$ (cor:vwz) rests on T5, which is proved for the free
fermion only.  None of these notes changes a status.

\section{Numerics: independent recomputations and the rerun}\label{sec:num}
Interpreter: \texttt{\$V/bin/python} (numpy 2.5.3, scipy 1.18.1, mpmath 1.4.1), $V$ the venv of the brief.  Commands
(repository root), each output saved next to its script:
\begin{itemize}[leftmargin=1.5em]
\item \texttt{\$\{PYTHON:-python3\} rigor/ref\_p6\_4\_kernel.py > rigor/ref\_p6\_4\_kernel.out} (47\,s): items V, P, R, K, B, E.
\item \texttt{\$\{PYTHON:-python3\} rigor/ref\_p6\_4\_composite.py > rigor/ref\_p6\_4\_composite.out} (1\,s): items L, N, W.
\item \texttt{PYTHON=\$V/bin/python \$V/bin/python rigor/ref\_p6\_4\_rerun\_g1b.py > rigor/ref\_p6\_4\_g1b\_kernel\_rerun.out}
 (under 2\,min).  \texttt{g1b\_kernel.py} writes \texttt{rigor/g1b\_kernel.out} \emph{itself} (its l.~206), so the brief's literal
 command would have overwritten the reviewed file; the wrapper executes the unchanged script by \texttt{runpy} and
 redirects that one write-open to \texttt{rigor/ref\_p6\_4\_g1b\_kernel\_rerun\_selfwrite.out}.  Result: both the stdout capture
 and the redirected file are \emph{byte-identical} to \texttt{rigor/g1b\_kernel.out} (\texttt{diff} empty), whose sha256
 \texttt{7581c0e3\dots cbbf} is unchanged.
\end{itemize}
Tolerances were fixed in the script docstrings before the runs.  No operator is discretised and no regulator is used
anywhere below (continuum integrals, exact kernels, rational arithmetic), so the convergence/regulator rules of
AGENTS.md \S5 reduce to the truncation statements given; ``validate first'' is item V.
\begin{center}\small
\begin{tabularx}{\textwidth}{>{\raggedright\arraybackslash}p{0.06\textwidth} >{\raggedright\arraybackslash}X >{\raggedright\arraybackslash}p{0.30\textwidth}}\toprule
item & quantity, method & result (tolerance)\\\midrule
V & $\int_0^\infty\cos(q\Delta)/(1+q^2)=\frac\pi2e^{-\Delta}$ ($\Delta=0.5,3$); $\int_0^\infty\tanh(\pi q)/(q(1+q^2))=2$ ($Q=20,40,80$);
 $\int_0^\infty\sin(pz)/\sinh\frac z2=\pi\tanh\pi p$ & max dev.\ $8.9\cdot10^{-16}$ ($10^{-12}$) PASS\\
P & $\Gh(\Delta)/\Gh(0)$ at 9 $\Delta$'s: Gauss--Legendre + exact Ci/$E_1$ tail vs closed form vs series & $7.2\cdot10^{-16}$
 ($10^{-8}$); series $5.0\cdot10^{-49}$ ($10^{-12}$) PASS\\
R & $\rho(q)$ from the $p$-integral, $q=0.1,0.5,1,2,5$ & rel.\ $2.7\cdot10^{-51}$ ($10^{-10}$) PASS\\
K & $\hat d_0$ by direct 2-D quadrature of eq.~(ddot), 6 points & $7.2\cdot10^{-18}$ ($10^{-9}$) PASS; sign flip excluded\\
B & T4(d) bounds on $\Delta=0.01,\dots,20$ (2000 pts) & 0 violations; sharpness to $8.5\cdot10^{-19}$ at $\Delta=40$\\
E & sign, monotonicity, $\Gh''$ log-singularity & 0/0 failures; increment $\to-\ln10/(24\pi^2)$, 7 digits\\
L & lem:Dtot, $M=2,3$, 8 points each, complex step & rel.\ $2.0\cdot10^{-50}$ ($10^{-12}$) PASS; wrong frame 133\,\%\\
N & lem:nogroup masses, exact & $-\frac32=-2t\kappa_2$ vs $-\frac{31}{14}$\\
W & VWZ corner data, exact, 4 $(a,b)$ & all equal to cor:vwz\\\bottomrule
\end{tabularx}
\end{center}
\noindent Values (three methods, CF, SR and QG, agree to $7\cdot10^{-16}$): $\Gh(\Delta)/\Gh(0)=0.898386580099973$,
$0.745615141121157$, $0.600745363582704$, $0.476960337995355$, $0.294523053439204$, $0.179784303566026$,
$0.109299033378915$, $0.066349836879668$, $0.024419213316489$ at $\Delta=0.5,1,1.5,2,3,4,5,6,8$; $\Gh(0)/f_2=1$.
MINOR on \texttt{g1b\_kernel.py}: its docstring and SUMMARY line call method B ``QAWF (semi-infinite)'', but production B
calls \texttt{quad(\dots, 0, UCUT, weight='cos')} (QAWO on $[0,10^5]$, its l.~81), as \doc\ \S9 correctly says; only V1d uses QAWF.

\section{Second pass (REF-P6-4b, 2026-10-05)}\label{sec:pass2}
Object: \doc\ as revised by G1b (29 pp.; log: no \verb|^!|, no undefined reference; theorem numbers unchanged
per its \texttt{.aux}), its new \S11 (\texttt{sec:corrREF}, l.~1252, items R1--R6), \texttt{g1b\_kernel.py/.out},
\texttt{findings\_entries.tex}, the cards.  Only the coordinator's five questions; actor \texttt{REF-P6-4b}.
\paragraph{(1) thm:T3 as repaired (l.~688).}
\lstep{1}{1}{The only step that uses RR \texttt{lem:groupbound} is $\langle1\rangle5$; under its new assumption
 $s\zeta_{\max}\norm{[\Pm,K]}_2<1$ it is valid verbatim.}
\lby{per factor $\sigma=s\zeta_k\le s\zeta_{\max}$: $\norm{G_{\pm\sigma}}_2\le\sigma\norm{[\Pm,K]}_2<1$ (Duhamel, RR
 $\langle1\rangle2$, every $s$); then both blocks of $e^{\pm\sigma K}$ have norm $<1$, $\Pp e^{\pm\sigma K}\Pp$ is
 invertible, index $0$ (RR C3(ii)), the two block norms agree, $\norm{G}_2^2=\norm{\Pm U\Pp}_2^2+\norm{\Pp U\Pm}_2^2=
 2\norm{\Pm U\Pp}_2^2$, hence $\norm{\Pm e^{\pm\sigma K}\Pp}_2\le\sigma\norm{[\Pm,K]}_2/\sqrt2=\sigma\norm{\Pm K\Pp}_2$;
 conjugation by $W_{d_k}$ and the decomposition $\Pm AB\Pp=\Pm A\Pp B\Pp+\Pm A\Pm B\Pp$ give $sZ\norm{\Pm K\Pp}_2$.
 Box [K4] now states the range}
\lstep{1}{2}{$\langle1\rangle7$ is valid verbatim.}
\lby{the hypotheses are $2\hat f_2<(sZ)^{-2}$ and $2\hat f_2<(2s^2\zeta_{\max}^2)^{-1}$, i.e.\ $2\hat f_2<m$; since
 $2\hat f_2$ is the infimum (PA eq.~(infimum)), $h'$ with $2\hat f_2\le\norm{\Pm K\Pp}_2^2<m$ arbitrarily close to $2\hat f_2$
 exist; for them $s^2Z^2\norm{\Pm K\Pp}_2^2<1$ and $\sqrt2s\zeta_{\max}\norm{\Pm K\Pp}_2<1$; continuity of
 $-\frac12\log(1-x)$}
\lstep{1}{3}{MINOR (wording): $\langle1\rangle6$ reads ``If $s^2Z^2\norm{\Pm K\Pp}_2^2<1$: $\Phi_{\rm tot}\le\dots$'' and
 cites $\langle1\rangle5$, whose new assumption it does not carry; read alone it is the old unrestricted claim for a
 fixed $h'$ (e.g.\ $M=1$, $\frac12\le s^2\zeta^2\norm{\Pm K\Pp}_2^2<1$).}
\lby{l.~725; harmless, because $\langle1\rangle7$ invokes $\langle1\rangle6$ only for $h'$ satisfying both conditions and
 says so (``$\langle1\rangle5$ applies'').  Fix: add ``and $s\zeta_{\max}\norm{[\Pm,K]}_2<1$'' to $\langle1\rangle6$}
\lstep{1}{4}{No other step needs index zero.}
\lby{$\langle1\rangle2$ (Dyson cocycle, both blocks, every $s\ge0$), $\langle1\rangle3$, and $\langle1\rangle4$ (PA thm:det
 for the product, with both block norms $<1$ supplied by $\langle1\rangle5$, so $\Pp U(s)\Pp$ is invertible)}
\begin{verdictok}
thm:T3: STANDS in the stated range $2\hat f_2s^2Z^2<1$, $s\zeta_{\max}<\pi\sqrt3$.  MINOR wording in $\langle1\rangle6$.
\end{verdictok}
\paragraph{(2) lem:onegroup(i) (l.~955) and thm:T5 $\langle1\rangle4$.}
\lstep{1}{5}{The side conditions hold for $s\le s_\eps$ uniformly on compact $\mathcal Z$ and in $y_1<y_2$.}
\lby{$h'$, hence $k$, $\norm k_2$, $\eta$, $\theta$, is fixed by $\delta$ alone ($\langle1\rangle3$); with
 $\bar z:=\max_{\mathcal Z}(\zeta_1+\zeta_2)$, every $s<(\sqrt2\,\bar z\norm k_2)^{-1}$ gives $s(\zeta_1+\zeta_2)\norm k_2<1$,
 $\sqrt2s\max_j\zeta_j\norm k_2<1$, and, by T3 $\langle1\rangle5$, $\norm V_2\le s(\zeta_1+\zeta_2)\norm k_2<2^{-1/2}$; none
 involves $y_1,y_2$; lem:onegroup(i) in $\langle2\rangle1$ needs $\sqrt2s\zeta_j\norm k_2<1$, the third condition}
\lstep{1}{6}{MINOR (wording, as in $\langle1\rangle3$): T5 $\langle1\rangle5$ is stated with no condition on $s$ but uses
 lem:onegroup(i) in $\langle2\rangle1$; it holds under the side conditions of $\langle1\rangle4$, which is how
 $\langle1\rangle6$ uses it (``choose $\delta$, then $s_\eps$'').}
\lby{T5 proof, read}
\begin{verdictok}
lem:onegroup and thm:T5: STAND; the T5 Statement needs nothing beyond $s\le s_\eps$.
\end{verdictok}
\paragraph{(3) thm:T4(e), new steps $\langle1\rangle4$--$\langle1\rangle5$.}
\lstep{1}{7}{The $C^1$/not-$C^2$ proof is complete, and (e) now states it.}
\lby{$\langle1\rangle4$: the (cf) series and its termwise derivative converge uniformly on $[0,\infty)$ (M-test with
 $\sum1/(a_n\abs{1-a_n^2}),\sum1/\abs{1-a_n^2}<\infty$), so $\Gh\in C^1[0,\infty)$ with right derivative
 $-\frac1{24\pi^2}\sum\frac1{1-a_n^2}=0$ at $0$, and evenness gives $C^1(\mathbb R)$, $\Gh'(0)=0$; $\langle1\rangle5$: if
 $\Gh\in C^2$ near $0$, $(\Gh(\Delta)-\Gh(0))/\Delta^2\to\frac12\Gh''(0)$; the expansions of $\cosh$, $\sinh^2$ and
 $\log\coth z=-\log z+\frac{z^2}3+O(z^4)$ are right and give $\frac18-\frac14\log\frac4{\abs\Delta}+o(1)\to-\infty$; first-pass
 numerics agree (\S\ref{sec:num}, item E)}
\paragraph{(4) \S11, R1--R6.}
\lstep{1}{8}{R1--R6 state my findings and the repairs accurately.}
\lby{R1 reproduces my locators, numbers and both uncovered examples and option (a); R2--R4 match
 \S\S\ref{sec:i10},\ref{sec:i12}; R5 verified: the current \texttt{g1b\_kernel.out} differs from the original only in its
 SUMMARY line (QAWF$\to$QAWO), and my guarded rerun of the corrected script
 (\texttt{rigor/ref\_p6\_4\_b\_rerun\_g1b.py} $\to$ \texttt{rigor/ref\_p6\_4\_b\_g1b\_kernel\_rerun.out}) reproduces it byte for
 byte (sha256 \texttt{e72d7ecd\dots}, unchanged); R6: the third G1b findings entry now carries ``for
 $s\max_k\zeta_k<\pi\sqrt3$ (and $2s^2(\sum\zeta_k)^2/(12\pi^2)<1$)'' with a dated note, and card
 \texttt{network-second-order-law} has a new note of 2026-10-05 with the corrected range.  ``Nothing else altered'' cannot
 be checked against a stored copy (the file is untracked); every spot check (T3, [K4], lem:onegroup, T5, T4(e),
 rem:kernelshape, Summary table l.~1222) agrees.  Not in R1--R6, and not required: my remark on the wording of G1b's note
 on \texttt{vwz-protocol-ordering} (\S\ref{sec:i16}); the abstract (l.~31) still says the coincident-corner bound ``is
 proved'' without its range}
\paragraph{(5) Card \texttt{network-global-bound-coincident}.}
\lstep{1}{9}{Statement (1) is thm:T3 with both conditions and names the excess range as open (card
 \texttt{group-index-zero-large-s}); \texttt{proved} is the weakest status the evidence supports for it; \texttt{depends}
 (\texttt{rate-of-remainder}, \texttt{thm-a-quadratic-law}, \texttt{network-law-free-fermion},
 \texttt{corner-kernel-closed-form}) is complete (G1a enters through \texttt{network-law-free-fermion}); locators
 \texttt{\#thm:T3} (l.~688), \texttt{\#lem:nogroup} (l.~664) and the How-to-verify's \texttt{sec:corrREF}/R1 (l.~1252,
 1257) resolve.}
\lby{\texttt{kb show}, \texttt{grep -n}}
\begin{verdictok}
Second pass: \texttt{network-global-bound-coincident} PASSES (\texttt{kb verdict} by REF-P6-4b), with the MINOR wording
items $\langle1\rangle3$ and $\langle1\rangle6$ recorded in the verdict note.
\end{verdictok}

\end{document}
